% Options for packages loaded elsewhere
\PassOptionsToPackage{unicode}{hyperref}
\PassOptionsToPackage{hyphens}{url}
\documentclass[
]{article}
\usepackage{xcolor}
\usepackage[margin=1in]{geometry}
\usepackage{amsmath,amssymb}
\setcounter{secnumdepth}{-\maxdimen} % remove section numbering
\usepackage{iftex}
\ifPDFTeX
  \usepackage[T1]{fontenc}
  \usepackage[utf8]{inputenc}
  \usepackage{textcomp} % provide euro and other symbols
\else % if luatex or xetex
  \usepackage{unicode-math} % this also loads fontspec
  \defaultfontfeatures{Scale=MatchLowercase}
  \defaultfontfeatures[\rmfamily]{Ligatures=TeX,Scale=1}
\fi
\usepackage{lmodern}
\ifPDFTeX\else
  % xetex/luatex font selection
\fi
% Use upquote if available, for straight quotes in verbatim environments
\IfFileExists{upquote.sty}{\usepackage{upquote}}{}
\IfFileExists{microtype.sty}{% use microtype if available
  \usepackage[]{microtype}
  \UseMicrotypeSet[protrusion]{basicmath} % disable protrusion for tt fonts
}{}
\makeatletter
\@ifundefined{KOMAClassName}{% if non-KOMA class
  \IfFileExists{parskip.sty}{%
    \usepackage{parskip}
  }{% else
    \setlength{\parindent}{0pt}
    \setlength{\parskip}{6pt plus 2pt minus 1pt}}
}{% if KOMA class
  \KOMAoptions{parskip=half}}
\makeatother
\usepackage{color}
\usepackage{fancyvrb}
\newcommand{\VerbBar}{|}
\newcommand{\VERB}{\Verb[commandchars=\\\{\}]}
\DefineVerbatimEnvironment{Highlighting}{Verbatim}{commandchars=\\\{\}}
% Add ',fontsize=\small' for more characters per line
\usepackage{framed}
\definecolor{shadecolor}{RGB}{248,248,248}
\newenvironment{Shaded}{\begin{snugshade}}{\end{snugshade}}
\newcommand{\AlertTok}[1]{\textcolor[rgb]{0.94,0.16,0.16}{#1}}
\newcommand{\AnnotationTok}[1]{\textcolor[rgb]{0.56,0.35,0.01}{\textbf{\textit{#1}}}}
\newcommand{\AttributeTok}[1]{\textcolor[rgb]{0.13,0.29,0.53}{#1}}
\newcommand{\BaseNTok}[1]{\textcolor[rgb]{0.00,0.00,0.81}{#1}}
\newcommand{\BuiltInTok}[1]{#1}
\newcommand{\CharTok}[1]{\textcolor[rgb]{0.31,0.60,0.02}{#1}}
\newcommand{\CommentTok}[1]{\textcolor[rgb]{0.56,0.35,0.01}{\textit{#1}}}
\newcommand{\CommentVarTok}[1]{\textcolor[rgb]{0.56,0.35,0.01}{\textbf{\textit{#1}}}}
\newcommand{\ConstantTok}[1]{\textcolor[rgb]{0.56,0.35,0.01}{#1}}
\newcommand{\ControlFlowTok}[1]{\textcolor[rgb]{0.13,0.29,0.53}{\textbf{#1}}}
\newcommand{\DataTypeTok}[1]{\textcolor[rgb]{0.13,0.29,0.53}{#1}}
\newcommand{\DecValTok}[1]{\textcolor[rgb]{0.00,0.00,0.81}{#1}}
\newcommand{\DocumentationTok}[1]{\textcolor[rgb]{0.56,0.35,0.01}{\textbf{\textit{#1}}}}
\newcommand{\ErrorTok}[1]{\textcolor[rgb]{0.64,0.00,0.00}{\textbf{#1}}}
\newcommand{\ExtensionTok}[1]{#1}
\newcommand{\FloatTok}[1]{\textcolor[rgb]{0.00,0.00,0.81}{#1}}
\newcommand{\FunctionTok}[1]{\textcolor[rgb]{0.13,0.29,0.53}{\textbf{#1}}}
\newcommand{\ImportTok}[1]{#1}
\newcommand{\InformationTok}[1]{\textcolor[rgb]{0.56,0.35,0.01}{\textbf{\textit{#1}}}}
\newcommand{\KeywordTok}[1]{\textcolor[rgb]{0.13,0.29,0.53}{\textbf{#1}}}
\newcommand{\NormalTok}[1]{#1}
\newcommand{\OperatorTok}[1]{\textcolor[rgb]{0.81,0.36,0.00}{\textbf{#1}}}
\newcommand{\OtherTok}[1]{\textcolor[rgb]{0.56,0.35,0.01}{#1}}
\newcommand{\PreprocessorTok}[1]{\textcolor[rgb]{0.56,0.35,0.01}{\textit{#1}}}
\newcommand{\RegionMarkerTok}[1]{#1}
\newcommand{\SpecialCharTok}[1]{\textcolor[rgb]{0.81,0.36,0.00}{\textbf{#1}}}
\newcommand{\SpecialStringTok}[1]{\textcolor[rgb]{0.31,0.60,0.02}{#1}}
\newcommand{\StringTok}[1]{\textcolor[rgb]{0.31,0.60,0.02}{#1}}
\newcommand{\VariableTok}[1]{\textcolor[rgb]{0.00,0.00,0.00}{#1}}
\newcommand{\VerbatimStringTok}[1]{\textcolor[rgb]{0.31,0.60,0.02}{#1}}
\newcommand{\WarningTok}[1]{\textcolor[rgb]{0.56,0.35,0.01}{\textbf{\textit{#1}}}}
\usepackage{graphicx}
\makeatletter
\newsavebox\pandoc@box
\newcommand*\pandocbounded[1]{% scales image to fit in text height/width
  \sbox\pandoc@box{#1}%
  \Gscale@div\@tempa{\textheight}{\dimexpr\ht\pandoc@box+\dp\pandoc@box\relax}%
  \Gscale@div\@tempb{\linewidth}{\wd\pandoc@box}%
  \ifdim\@tempb\p@<\@tempa\p@\let\@tempa\@tempb\fi% select the smaller of both
  \ifdim\@tempa\p@<\p@\scalebox{\@tempa}{\usebox\pandoc@box}%
  \else\usebox{\pandoc@box}%
  \fi%
}
% Set default figure placement to htbp
\def\fps@figure{htbp}
\makeatother
\setlength{\emergencystretch}{3em} % prevent overfull lines
\providecommand{\tightlist}{%
  \setlength{\itemsep}{0pt}\setlength{\parskip}{0pt}}
\usepackage[spanish]{babel}
\usepackage{amsmath}
\usepackage{fontspec}
\usepackage{unicode-math}
\usepackage{graphicx}
\usepackage{adjustbox}
\usepackage{tikz}
\usepackage{pgfplots}
\usetikzlibrary{3d,babel}
\usepackage{graphicx}
\usepackage{float}
\usepackage{tikz}
\usepackage{xcolor}
\usepackage{amsmath}
\usepackage{pgfplots}
\usepackage{bookmark}
\IfFileExists{xurl.sty}{\usepackage{xurl}}{} % add URL line breaks if available
\urlstyle{same}
\hypersetup{
  hidelinks,
  pdfcreator={LaTeX via pandoc}}

\author{}
\date{\vspace{-2.5em}}

\begin{document}

\begin{Shaded}
\begin{Highlighting}[]
\FunctionTok{options}\NormalTok{(}\AttributeTok{OutDec =} \StringTok{"."}\NormalTok{)  }\CommentTok{\# Asegurar punto decimal en este chunk}

\CommentTok{\# Función para generar datos aleatorios del problema}
\NormalTok{generar\_datos }\OtherTok{\textless{}{-}} \ControlFlowTok{function}\NormalTok{() \{}
  \CommentTok{\# Generar pendiente aleatoria (entre {-}3 y 3, evitando 0)}
\NormalTok{  pendientes\_posibles }\OtherTok{\textless{}{-}} \FunctionTok{c}\NormalTok{(}\SpecialCharTok{{-}}\DecValTok{3}\NormalTok{, }\SpecialCharTok{{-}}\FloatTok{2.5}\NormalTok{, }\SpecialCharTok{{-}}\DecValTok{2}\NormalTok{, }\SpecialCharTok{{-}}\FloatTok{1.5}\NormalTok{, }\SpecialCharTok{{-}}\DecValTok{1}\NormalTok{, }\SpecialCharTok{{-}}\FloatTok{0.5}\NormalTok{, }\FloatTok{0.5}\NormalTok{, }\DecValTok{1}\NormalTok{, }\FloatTok{1.5}\NormalTok{, }\DecValTok{2}\NormalTok{, }\FloatTok{2.5}\NormalTok{, }\DecValTok{3}\NormalTok{)}
\NormalTok{  m }\OtherTok{\textless{}{-}} \FunctionTok{sample}\NormalTok{(pendientes\_posibles, }\DecValTok{1}\NormalTok{)}
  
  \CommentTok{\# Generar intercepto aleatorio (entre {-}4 y 4)}
\NormalTok{  b }\OtherTok{\textless{}{-}} \FunctionTok{sample}\NormalTok{(}\FunctionTok{seq}\NormalTok{(}\SpecialCharTok{{-}}\DecValTok{4}\NormalTok{, }\DecValTok{4}\NormalTok{, }\AttributeTok{by =} \FloatTok{0.5}\NormalTok{), }\DecValTok{1}\NormalTok{)}
  
  \CommentTok{\# Formatear pendiente e intercepto para mostrar}
\NormalTok{  m\_formateado }\OtherTok{\textless{}{-}} \FunctionTok{ifelse}\NormalTok{(m }\SpecialCharTok{==} \FunctionTok{as.integer}\NormalTok{(m), }\FunctionTok{as.character}\NormalTok{(m), }
                        \FunctionTok{ifelse}\NormalTok{(}\FunctionTok{abs}\NormalTok{(m) }\SpecialCharTok{==} \FloatTok{0.5}\NormalTok{, }\FunctionTok{ifelse}\NormalTok{(m }\SpecialCharTok{\textless{}} \DecValTok{0}\NormalTok{, }\StringTok{"{-}}\SpecialCharTok{\textbackslash{}\textbackslash{}}\StringTok{frac\{1\}\{2\}"}\NormalTok{, }\StringTok{"}\SpecialCharTok{\textbackslash{}\textbackslash{}}\StringTok{frac\{1\}\{2\}"}\NormalTok{),}
                               \FunctionTok{ifelse}\NormalTok{(}\FunctionTok{abs}\NormalTok{(m) }\SpecialCharTok{==} \FloatTok{1.5}\NormalTok{, }\FunctionTok{ifelse}\NormalTok{(m }\SpecialCharTok{\textless{}} \DecValTok{0}\NormalTok{, }\StringTok{"{-}}\SpecialCharTok{\textbackslash{}\textbackslash{}}\StringTok{frac\{3\}\{2\}"}\NormalTok{, }\StringTok{"}\SpecialCharTok{\textbackslash{}\textbackslash{}}\StringTok{frac\{3\}\{2\}"}\NormalTok{),}
                                      \FunctionTok{ifelse}\NormalTok{(}\FunctionTok{abs}\NormalTok{(m) }\SpecialCharTok{==} \FloatTok{2.5}\NormalTok{, }\FunctionTok{ifelse}\NormalTok{(m }\SpecialCharTok{\textless{}} \DecValTok{0}\NormalTok{, }\StringTok{"{-}}\SpecialCharTok{\textbackslash{}\textbackslash{}}\StringTok{frac\{5\}\{2\}"}\NormalTok{, }\StringTok{"}\SpecialCharTok{\textbackslash{}\textbackslash{}}\StringTok{frac\{5\}\{2\}"}\NormalTok{),}
                                             \FunctionTok{sprintf}\NormalTok{(}\StringTok{"\%.1f"}\NormalTok{, m)))))}
  
\NormalTok{  b\_formateado }\OtherTok{\textless{}{-}} \FunctionTok{ifelse}\NormalTok{(b }\SpecialCharTok{==} \FunctionTok{as.integer}\NormalTok{(b), }\FunctionTok{as.character}\NormalTok{(b), }\FunctionTok{sprintf}\NormalTok{(}\StringTok{"\%.1f"}\NormalTok{, b))}
  
  \CommentTok{\# Generar ecuación de la recta}
  \ControlFlowTok{if}\NormalTok{ (b }\SpecialCharTok{==} \DecValTok{0}\NormalTok{) \{}
\NormalTok{    ecuacion }\OtherTok{\textless{}{-}} \FunctionTok{paste0}\NormalTok{(}\StringTok{"y = "}\NormalTok{, m\_formateado, }\StringTok{"x"}\NormalTok{)}
\NormalTok{  \} }\ControlFlowTok{else} \ControlFlowTok{if}\NormalTok{ (b }\SpecialCharTok{\textgreater{}} \DecValTok{0}\NormalTok{) \{}
\NormalTok{    ecuacion }\OtherTok{\textless{}{-}} \FunctionTok{paste0}\NormalTok{(}\StringTok{"y = "}\NormalTok{, m\_formateado, }\StringTok{"x + "}\NormalTok{, b\_formateado)}
\NormalTok{  \} }\ControlFlowTok{else}\NormalTok{ \{}
    \CommentTok{\# Usar abs() sobre el valor numérico, no sobre el string formateado}
\NormalTok{    b\_abs\_formateado }\OtherTok{\textless{}{-}} \FunctionTok{ifelse}\NormalTok{(}\FunctionTok{abs}\NormalTok{(b) }\SpecialCharTok{==} \FunctionTok{as.integer}\NormalTok{(}\FunctionTok{abs}\NormalTok{(b)), }
                               \FunctionTok{as.character}\NormalTok{(}\FunctionTok{abs}\NormalTok{(b)), }
                               \FunctionTok{sprintf}\NormalTok{(}\StringTok{"\%.1f"}\NormalTok{, }\FunctionTok{abs}\NormalTok{(b)))}
\NormalTok{    ecuacion }\OtherTok{\textless{}{-}} \FunctionTok{paste0}\NormalTok{(}\StringTok{"y = "}\NormalTok{, m\_formateado, }\StringTok{"x {-} "}\NormalTok{, b\_abs\_formateado)}
\NormalTok{  \}}
  
  \CommentTok{\# Generar opciones de respuesta (distractores pedagógicamente efectivos)}
  \CommentTok{\# Opción correcta: la ecuación completa}
\NormalTok{  opcion\_correcta }\OtherTok{\textless{}{-}}\NormalTok{ ecuacion}
  
  \CommentTok{\# Distractor 1: Pendiente correcta pero intercepto incorrecto}
\NormalTok{  b\_distractor1 }\OtherTok{\textless{}{-}} \FunctionTok{sample}\NormalTok{(}\FunctionTok{c}\NormalTok{(b }\SpecialCharTok{+} \DecValTok{1}\NormalTok{, b }\SpecialCharTok{{-}} \DecValTok{1}\NormalTok{, b }\SpecialCharTok{+} \DecValTok{2}\NormalTok{, b }\SpecialCharTok{{-}} \DecValTok{2}\NormalTok{), }\DecValTok{1}\NormalTok{)}
\NormalTok{  b\_dist1\_formateado }\OtherTok{\textless{}{-}} \FunctionTok{ifelse}\NormalTok{(b\_distractor1 }\SpecialCharTok{==} \FunctionTok{as.integer}\NormalTok{(b\_distractor1), }
                               \FunctionTok{as.character}\NormalTok{(b\_distractor1), }
                               \FunctionTok{sprintf}\NormalTok{(}\StringTok{"\%.1f"}\NormalTok{, b\_distractor1))}
  \ControlFlowTok{if}\NormalTok{ (b\_distractor1 }\SpecialCharTok{==} \DecValTok{0}\NormalTok{) \{}
\NormalTok{    distractor1 }\OtherTok{\textless{}{-}} \FunctionTok{paste0}\NormalTok{(}\StringTok{"y = "}\NormalTok{, m\_formateado, }\StringTok{"x"}\NormalTok{)}
\NormalTok{  \} }\ControlFlowTok{else} \ControlFlowTok{if}\NormalTok{ (b\_distractor1 }\SpecialCharTok{\textgreater{}} \DecValTok{0}\NormalTok{) \{}
\NormalTok{    distractor1 }\OtherTok{\textless{}{-}} \FunctionTok{paste0}\NormalTok{(}\StringTok{"y = "}\NormalTok{, m\_formateado, }\StringTok{"x + "}\NormalTok{, b\_dist1\_formateado)}
\NormalTok{  \} }\ControlFlowTok{else}\NormalTok{ \{}
\NormalTok{    b\_dist1\_abs\_formateado }\OtherTok{\textless{}{-}} \FunctionTok{ifelse}\NormalTok{(}\FunctionTok{abs}\NormalTok{(b\_distractor1) }\SpecialCharTok{==} \FunctionTok{as.integer}\NormalTok{(}\FunctionTok{abs}\NormalTok{(b\_distractor1)), }
                                     \FunctionTok{as.character}\NormalTok{(}\FunctionTok{abs}\NormalTok{(b\_distractor1)), }
                                     \FunctionTok{sprintf}\NormalTok{(}\StringTok{"\%.1f"}\NormalTok{, }\FunctionTok{abs}\NormalTok{(b\_distractor1)))}
\NormalTok{    distractor1 }\OtherTok{\textless{}{-}} \FunctionTok{paste0}\NormalTok{(}\StringTok{"y = "}\NormalTok{, m\_formateado, }\StringTok{"x {-} "}\NormalTok{, b\_dist1\_abs\_formateado)}
\NormalTok{  \}}
  
  \CommentTok{\# Distractor 2: Intercepto correcto pero pendiente incorrecta}
\NormalTok{  m\_distractor2 }\OtherTok{\textless{}{-}} \FunctionTok{sample}\NormalTok{(pendientes\_posibles[pendientes\_posibles }\SpecialCharTok{!=}\NormalTok{ m], }\DecValTok{1}\NormalTok{)}
\NormalTok{  m\_dist2\_formateado }\OtherTok{\textless{}{-}} \FunctionTok{ifelse}\NormalTok{(m\_distractor2 }\SpecialCharTok{==} \FunctionTok{as.integer}\NormalTok{(m\_distractor2), }
                               \FunctionTok{as.character}\NormalTok{(m\_distractor2),}
                               \FunctionTok{ifelse}\NormalTok{(}\FunctionTok{abs}\NormalTok{(m\_distractor2) }\SpecialCharTok{==} \FloatTok{0.5}\NormalTok{, }
                                      \FunctionTok{ifelse}\NormalTok{(m\_distractor2 }\SpecialCharTok{\textless{}} \DecValTok{0}\NormalTok{, }\StringTok{"{-}}\SpecialCharTok{\textbackslash{}\textbackslash{}}\StringTok{frac\{1\}\{2\}"}\NormalTok{, }\StringTok{"}\SpecialCharTok{\textbackslash{}\textbackslash{}}\StringTok{frac\{1\}\{2\}"}\NormalTok{),}
                                      \FunctionTok{ifelse}\NormalTok{(}\FunctionTok{abs}\NormalTok{(m\_distractor2) }\SpecialCharTok{==} \FloatTok{1.5}\NormalTok{, }
                                             \FunctionTok{ifelse}\NormalTok{(m\_distractor2 }\SpecialCharTok{\textless{}} \DecValTok{0}\NormalTok{, }\StringTok{"{-}}\SpecialCharTok{\textbackslash{}\textbackslash{}}\StringTok{frac\{3\}\{2\}"}\NormalTok{, }\StringTok{"}\SpecialCharTok{\textbackslash{}\textbackslash{}}\StringTok{frac\{3\}\{2\}"}\NormalTok{),}
                                             \FunctionTok{ifelse}\NormalTok{(}\FunctionTok{abs}\NormalTok{(m\_distractor2) }\SpecialCharTok{==} \FloatTok{2.5}\NormalTok{, }
                                                    \FunctionTok{ifelse}\NormalTok{(m\_distractor2 }\SpecialCharTok{\textless{}} \DecValTok{0}\NormalTok{, }\StringTok{"{-}}\SpecialCharTok{\textbackslash{}\textbackslash{}}\StringTok{frac\{5\}\{2\}"}\NormalTok{, }\StringTok{"}\SpecialCharTok{\textbackslash{}\textbackslash{}}\StringTok{frac\{5\}\{2\}"}\NormalTok{),}
                                                    \FunctionTok{sprintf}\NormalTok{(}\StringTok{"\%.1f"}\NormalTok{, m\_distractor2)))))}
  \ControlFlowTok{if}\NormalTok{ (b }\SpecialCharTok{==} \DecValTok{0}\NormalTok{) \{}
\NormalTok{    distractor2 }\OtherTok{\textless{}{-}} \FunctionTok{paste0}\NormalTok{(}\StringTok{"y = "}\NormalTok{, m\_dist2\_formateado, }\StringTok{"x"}\NormalTok{)}
\NormalTok{  \} }\ControlFlowTok{else} \ControlFlowTok{if}\NormalTok{ (b }\SpecialCharTok{\textgreater{}} \DecValTok{0}\NormalTok{) \{}
\NormalTok{    distractor2 }\OtherTok{\textless{}{-}} \FunctionTok{paste0}\NormalTok{(}\StringTok{"y = "}\NormalTok{, m\_dist2\_formateado, }\StringTok{"x + "}\NormalTok{, b\_formateado)}
\NormalTok{  \} }\ControlFlowTok{else}\NormalTok{ \{}
\NormalTok{    b\_abs\_formateado }\OtherTok{\textless{}{-}} \FunctionTok{ifelse}\NormalTok{(}\FunctionTok{abs}\NormalTok{(b) }\SpecialCharTok{==} \FunctionTok{as.integer}\NormalTok{(}\FunctionTok{abs}\NormalTok{(b)), }
                               \FunctionTok{as.character}\NormalTok{(}\FunctionTok{abs}\NormalTok{(b)), }
                               \FunctionTok{sprintf}\NormalTok{(}\StringTok{"\%.1f"}\NormalTok{, }\FunctionTok{abs}\NormalTok{(b)))}
\NormalTok{    distractor2 }\OtherTok{\textless{}{-}} \FunctionTok{paste0}\NormalTok{(}\StringTok{"y = "}\NormalTok{, m\_dist2\_formateado, }\StringTok{"x {-} "}\NormalTok{, b\_abs\_formateado)}
\NormalTok{  \}}
  
  \CommentTok{\# Distractor 3: Ambos incorrectos pero con valores cercanos}
\NormalTok{  m\_distractor3 }\OtherTok{\textless{}{-}} \FunctionTok{sample}\NormalTok{(pendientes\_posibles[pendientes\_posibles }\SpecialCharTok{!=}\NormalTok{ m }\SpecialCharTok{\&}\NormalTok{ pendientes\_posibles }\SpecialCharTok{!=}\NormalTok{ m\_distractor2], }\DecValTok{1}\NormalTok{)}
\NormalTok{  b\_distractor3 }\OtherTok{\textless{}{-}} \FunctionTok{sample}\NormalTok{(}\FunctionTok{c}\NormalTok{(b }\SpecialCharTok{+} \DecValTok{1}\NormalTok{, b }\SpecialCharTok{{-}} \DecValTok{1}\NormalTok{, b }\SpecialCharTok{+} \DecValTok{2}\NormalTok{, b }\SpecialCharTok{{-}} \DecValTok{2}\NormalTok{), }\DecValTok{1}\NormalTok{)}
\NormalTok{  m\_dist3\_formateado }\OtherTok{\textless{}{-}} \FunctionTok{ifelse}\NormalTok{(m\_distractor3 }\SpecialCharTok{==} \FunctionTok{as.integer}\NormalTok{(m\_distractor3), }
                               \FunctionTok{as.character}\NormalTok{(m\_distractor3),}
                               \FunctionTok{ifelse}\NormalTok{(}\FunctionTok{abs}\NormalTok{(m\_distractor3) }\SpecialCharTok{==} \FloatTok{0.5}\NormalTok{, }
                                      \FunctionTok{ifelse}\NormalTok{(m\_distractor3 }\SpecialCharTok{\textless{}} \DecValTok{0}\NormalTok{, }\StringTok{"{-}}\SpecialCharTok{\textbackslash{}\textbackslash{}}\StringTok{frac\{1\}\{2\}"}\NormalTok{, }\StringTok{"}\SpecialCharTok{\textbackslash{}\textbackslash{}}\StringTok{frac\{1\}\{2\}"}\NormalTok{),}
                                      \FunctionTok{ifelse}\NormalTok{(}\FunctionTok{abs}\NormalTok{(m\_distractor3) }\SpecialCharTok{==} \FloatTok{1.5}\NormalTok{, }
                                             \FunctionTok{ifelse}\NormalTok{(m\_distractor3 }\SpecialCharTok{\textless{}} \DecValTok{0}\NormalTok{, }\StringTok{"{-}}\SpecialCharTok{\textbackslash{}\textbackslash{}}\StringTok{frac\{3\}\{2\}"}\NormalTok{, }\StringTok{"}\SpecialCharTok{\textbackslash{}\textbackslash{}}\StringTok{frac\{3\}\{2\}"}\NormalTok{),}
                                             \FunctionTok{ifelse}\NormalTok{(}\FunctionTok{abs}\NormalTok{(m\_distractor3) }\SpecialCharTok{==} \FloatTok{2.5}\NormalTok{, }
                                                    \FunctionTok{ifelse}\NormalTok{(m\_distractor3 }\SpecialCharTok{\textless{}} \DecValTok{0}\NormalTok{, }\StringTok{"{-}}\SpecialCharTok{\textbackslash{}\textbackslash{}}\StringTok{frac\{5\}\{2\}"}\NormalTok{, }\StringTok{"}\SpecialCharTok{\textbackslash{}\textbackslash{}}\StringTok{frac\{5\}\{2\}"}\NormalTok{),}
                                                    \FunctionTok{sprintf}\NormalTok{(}\StringTok{"\%.1f"}\NormalTok{, m\_distractor3)))))}
\NormalTok{  b\_dist3\_formateado }\OtherTok{\textless{}{-}} \FunctionTok{ifelse}\NormalTok{(b\_distractor3 }\SpecialCharTok{==} \FunctionTok{as.integer}\NormalTok{(b\_distractor3), }
                               \FunctionTok{as.character}\NormalTok{(b\_distractor3), }
                               \FunctionTok{sprintf}\NormalTok{(}\StringTok{"\%.1f"}\NormalTok{, b\_distractor3))}
  \ControlFlowTok{if}\NormalTok{ (b\_distractor3 }\SpecialCharTok{==} \DecValTok{0}\NormalTok{) \{}
\NormalTok{    distractor3 }\OtherTok{\textless{}{-}} \FunctionTok{paste0}\NormalTok{(}\StringTok{"y = "}\NormalTok{, m\_dist3\_formateado, }\StringTok{"x"}\NormalTok{)}
\NormalTok{  \} }\ControlFlowTok{else} \ControlFlowTok{if}\NormalTok{ (b\_distractor3 }\SpecialCharTok{\textgreater{}} \DecValTok{0}\NormalTok{) \{}
\NormalTok{    distractor3 }\OtherTok{\textless{}{-}} \FunctionTok{paste0}\NormalTok{(}\StringTok{"y = "}\NormalTok{, m\_dist3\_formateado, }\StringTok{"x + "}\NormalTok{, b\_dist3\_formateado)}
\NormalTok{  \} }\ControlFlowTok{else}\NormalTok{ \{}
\NormalTok{    b\_dist3\_abs\_formateado }\OtherTok{\textless{}{-}} \FunctionTok{ifelse}\NormalTok{(}\FunctionTok{abs}\NormalTok{(b\_distractor3) }\SpecialCharTok{==} \FunctionTok{as.integer}\NormalTok{(}\FunctionTok{abs}\NormalTok{(b\_distractor3)), }
                                     \FunctionTok{as.character}\NormalTok{(}\FunctionTok{abs}\NormalTok{(b\_distractor3)), }
                                     \FunctionTok{sprintf}\NormalTok{(}\StringTok{"\%.1f"}\NormalTok{, }\FunctionTok{abs}\NormalTok{(b\_distractor3)))}
\NormalTok{    distractor3 }\OtherTok{\textless{}{-}} \FunctionTok{paste0}\NormalTok{(}\StringTok{"y = "}\NormalTok{, m\_dist3\_formateado, }\StringTok{"x {-} "}\NormalTok{, b\_dist3\_abs\_formateado)}
\NormalTok{  \}}
  
  \CommentTok{\# Mezclar opciones}
\NormalTok{  todas\_opciones }\OtherTok{\textless{}{-}} \FunctionTok{c}\NormalTok{(opcion\_correcta, distractor1, distractor2, distractor3)}
\NormalTok{  orden\_opciones }\OtherTok{\textless{}{-}} \FunctionTok{sample}\NormalTok{(}\DecValTok{1}\SpecialCharTok{:}\DecValTok{4}\NormalTok{)}
\NormalTok{  opciones\_mezcladas }\OtherTok{\textless{}{-}}\NormalTok{ todas\_opciones[orden\_opciones]}
\NormalTok{  posicion\_correcta }\OtherTok{\textless{}{-}} \FunctionTok{which}\NormalTok{(orden\_opciones }\SpecialCharTok{==} \DecValTok{1}\NormalTok{)}
  
  \FunctionTok{return}\NormalTok{(}\FunctionTok{list}\NormalTok{(}
    \AttributeTok{m =}\NormalTok{ m,}
    \AttributeTok{b =}\NormalTok{ b,}
    \AttributeTok{m\_formateado =}\NormalTok{ m\_formateado,}
    \AttributeTok{b\_formateado =}\NormalTok{ b\_formateado,}
    \AttributeTok{ecuacion =}\NormalTok{ ecuacion,}
    \AttributeTok{opciones =}\NormalTok{ opciones\_mezcladas,}
    \AttributeTok{respuesta\_correcta =}\NormalTok{ opcion\_correcta,}
    \AttributeTok{posicion\_correcta =}\NormalTok{ posicion\_correcta}
\NormalTok{  ))}
\NormalTok{\}}

\CommentTok{\# Generar datos para este ejercicio}
\NormalTok{datos }\OtherTok{\textless{}{-}} \FunctionTok{generar\_datos}\NormalTok{()}
\end{Highlighting}
\end{Shaded}

\section{Question}\label{question}

En el siguiente gráfico se muestra una recta en el plano cartesiano:

\begin{figure}[H]

{\centering \includegraphics[width=0.5\linewidth]{recta_python} 

}

\end{figure}

¿Cuál es la ecuación de la recta representada en el gráfico?

\subsection{Answerlist}\label{answerlist}

\begin{itemize}
\tightlist
\item
  y = -1x - 3
\item
  y = -\frac{1}{2}x - 4
\item
  y = -\frac{3}{2}x - 5
\item
  y = -\frac{3}{2}x - 4
\end{itemize}

\section{Solution}\label{solution}

Para determinar la ecuación de la recta, necesitamos identificar su
pendiente (\(m\)) y su intercepto con el eje \(y\) (\(b\)).

La ecuación de una recta tiene la forma: \[y = mx + b\]

donde: - \(m\) es la pendiente de la recta - \(b\) es el intercepto con
el eje \(y\) (el valor de \(y\) cuando \(x = 0\))

\textbf{Análisis del gráfico:}

\begin{enumerate}
\def\labelenumi{\arabic{enumi}.}
\item
  \textbf{Intercepto con el eje \(y\)}: Observando el gráfico, la recta
  intersecta el eje \(y\) en el punto \((0, -4)\). Por lo tanto,
  \(b = -4\).
\item
  \textbf{Pendiente}: La pendiente se puede calcular usando dos puntos
  de la recta. Observando el gráfico:

  \begin{itemize}
  \tightlist
  \item
    La recta pasa por el punto \((0, -4)\)
  \item
    Para calcular la pendiente, podemos usar otro punto visible en la
    recta
  \end{itemize}

  La pendiente es \(m = -\frac{3}{2}\).
\end{enumerate}

\textbf{Ecuación de la recta:}

Sustituyendo los valores encontrados: \[y = y = -\frac{3}{2}x - 4\]

Por lo tanto, la ecuación de la recta es:
\textbf{\(y = -\frac{3}{2}x - 4\)}

\subsection{Answerlist}\label{answerlist-1}

\begin{itemize}
\tightlist
\item
  Incorrecto. No corresponde a la ecuación de la recta mostrada en el
  gráfico.
\item
  Incorrecto. Error en la identificación del intercepto o la pendiente.
\item
  Incorrecto. Error en la identificación del intercepto o la pendiente.
\item
  Correcto. Identifica correctamente la ecuación de la recta analizando
  la pendiente y el intercepto en el gráfico.
\end{itemize}

\section{Meta-information}\label{meta-information}

exname:
recta\_geometria\_analitica\_python\_interpretacion\_representacion
extype: schoice exsolution: 0001 exshuffle: TRUE exsection:
Geometría/Geometría Analítica/Ecuación de la recta

\end{document}
